\chapter{讨论证据边界}
\section{结果分析}
样例 A：误差下降可能来自门控、注意力或尾部损失；当前完整模型比较无法验证这些机制。

样例 B：误差变化与门控结构一致，局部轨迹为这一关联提供了证据锚点。

样例 C：移除门控后 RMSE 增大，关闭注意力后尾部误差增大，两项记录分别区分两个假设。

样例 D：当前结果不能区分形成该差异的机制，原因尚未确定。

样例 E：受控消融逐项移除门控模块后 RMSE 增大，因此该记录支持性能增益来自门控模块。

据此，下一节检查适用边界。
